% Глава: 21-docker-data
% Статус: наполнено (перый прогон).
\chapter{Данные в Docker: тома и монтирование}
\label{ch:21-docker-data}

\section{Проблема: данные в контейнере теряются}

Всё, что записано в файловую систему контейнера
(не в том, что было в образе), живёт в \emph{слое
контейнера}. Когда контейнер удаляется --- данные
исчезают. Для сохранения нужны \textbf{тома}
(volumes) и \textbf{монтирование} (mounts).

Тома размещаются на диске вашей ВМ на \fpath{pglab.nsu.ru};
учитывайте, что там 20--40~ГБ, и следите за местом.

\section{Три способа хранения}

\begin{center}
\begin{tabular}{lp{0.24\linewidth}lp{0.24\linewidth}}
\toprule
\textbf{Способ} & \textbf{Хранение} & \textbf{Управление} & \textbf{Кто управляет} \\
\midrule
named volume & в хранилище Docker & Docker & Docker \\
bind mount   & произвольный каталог хоста & вы & вы \\
tmpfs        & в RAM хоста & Docker & Docker \\
\bottomrule
\end{tabular}
\end{center}

\section{Именованные тома (named volumes)}

Создаются и управляются Docker:
\begin{minted}[bgcolor=gray!6]{bash}
docker volume create mydata          # создать
docker volume ls                     # список
docker volume inspect mydata         # подробно
docker volume rm mydata              # удалить, только пустой
\end{minted}

Подключить к контейнеру:
\begin{minted}[bgcolor=gray!6]{bash}
docker run -d -v mydata:/var/lib/app-data myimage
\end{minted}

\cmd{-v host:container} --- общий синтаксис:
что хост-путь (или имя тома), что путь внутри
контейнера.

\section{Bind mounts}

Монтирование \emph{произвольного} каталога хоста:
\begin{minted}[bgcolor=gray!6]{bash}
docker run -it -v /home/me/project:/app ubuntu ls /app
\end{minted}

Полезно для разработки: меняете файл на хосте ---
он сразу виден в контейнере.

\section{tmpfs}

\textbf{tmpfs} --- каталог в RAM; исчезает при
остановке контейнера:
\begin{minted}[bgcolor=gray!6]{bash}
docker run -it --tmpfs /tmp:size=100m ubuntu ls -la /tmp
\end{minted}

\section{Схема: хост, образ, контейнер, том}

\begin{figure}[htb]
  \centering
  \begin{tikzpicture}[node distance=1.4cm]
    \node[storage, minimum width=4cm] (vol) {named volume\\[0.2em] (on host disk)};
    \node[process, above right=1.2cm and 1cm of vol] (ctr) {container\\[0.2em] (ephemeral)};
    \node[data, above left=1.2cm and 1cm of vol] (img) {image (read-only)};
    \draw[flow, <->] (vol) -- node[right, font=\small]{mount} (ctr);
    \draw[flow, ->] (img) -- node[left, font=\small]{copy up} (ctr);
  \end{tikzpicture}
  \caption{Образ --- только чтение, слой контейнера --- временный,
    том --- постоянное хранилище.}
  \label{fig:docker-data}
\end{figure}

\section{Практика: база данных в Docker}

Типовой сценарий: образ с приложением + отдельный том
для данных:
\begin{minted}[bgcolor=gray!6]{bash}
docker volume create pg_data
docker run -d \\
  --name pg \\
  -v pg_data:/var/lib/postgresql/data \\
  -e POSTGRES_PASSWORD=secret \\
  postgres:16
\end{minted}

Даже если контейнер удалить, данные останутся в томе
\cmd{pg_data}.

\section{Типичные ошибки}
\begin{itemize}
  \item Данные «пропали» после \cmd{docker rm} --- не
    подключён том;
  \item Bind mount на каталог, которого нет на хосте ---
    Docker создаст его \emph{пустым};
  \item Права: приложение в контейнере работает от
    \cmd{root}, а на хосте каталог принадлежит вам ---
    возможны конфликты прав.
\end{itemize}

\section{Связь с остальным курсом}

Тома используются в \cmd{docker-compose}
(\nameref{ch:22-docker-compose}); ограничения на
ресурсы контейнеров --- \nameref{ch:23-docker-resources};
файловая система хоста --- \nameref{ch:04-filesystem}.

\section{Задачи для самопроверки}
\begin{enumerate}
  \item Почему данные в контейнере теряются после
    его удаления?
  \item Чем named volume отличается от bind mount?
  \item Как подключить том к контейнеру?
  \item Что такое tmpfs и когда его используют?
  \item Как сохранить данные PostgreSQL при удалении
    контейнера?
\end{enumerate}
